% Kapitel 2 — Budget: 4 Seiten — Schreiben: 14.–23.12.
% Eigenes, kurzes Kapitel (Beschluss 26.10.); nur, was der Leser zum Verstehen
% der Arbeit braucht; nur das Wichtige, Mut zur Lücke.

\chapter{Grundlagen}
\label{chap:grundlagen}

\section{Coding-Agent-Systeme: Funktionsweise und Aufgabenrahmen}% 2 S.
% Inhalt: bisher 2.1; wie ein Coding-Agent-System funktioniert,
% um den Ablauf der Arbeit zu verstehen. Text liegt vor, überarbeiten.
\label{sec:coding-agent-systeme}

\section{Token und Kontextfenster als relevante Währung}% 2 S., neu schreiben
% Inhalt: was die Zahl x Token genau bedeutet; teils aus bisherigem 2.2 und 2.4.
\label{sec:token-kontextfenster}

\section{Semantische und lexikalische Suche}% ca. 0,75 S., neu (Beschluss 30.10.)
% ===== WARUM DIESER ABSCHNITT =====
% Kapitel 4 benutzt Suchbegriffe aus dem Information Retrieval, die ein Leser
% ohne IR-Hintergrund nicht kennt. Bisher erklärt Kapitel 4 sie nebenbei im
% Fließtext (vor allem 4.4), das bläht 4.4 auf (Stand 30.10.: ca. 4 S. statt
% 2,5 S. Budget). Ziel: Begriffe hier EINMAL allgemein erklären, Kapitel 4
% verweist nur noch mit Abschnitt~\ref{sec:suche-grundlagen} zurück und
% beschreibt dort nur noch, WIE das Werkzeug sie einsetzt.
% Der erste Gebrauch liegt schon in 4.3 (sec:codegraph, Absatz "Beim Anlegen
% jedes Knotens erzeugt import.js dessen Einbettung ..."), nicht erst in 4.4.
% Kapitel 2 kommt vor 4.3, also muss hier alles stehen, was 4.3 und 4.4
% voraussetzen.
% Kapitel 3: in sec:graphbasiert prüfen, ob die dort besprochenen Arbeiten
% (RepoGraph, CodexGraph, KGCompass usw.) Einbettungen oder BM25 nutzen. Falls
% ja, dort ebenfalls hierher verweisen. Noch NICHT geprüft.
%
% ===== REGELN FÜR DEN TEXT =====
% - Nur allgemein erklären, nichts Werkzeugspezifisches (kein Modellname,
%   keine Schwellwerte, kein Neo4j). Das bleibt in Kapitel 4.
% - Jeder Begriff nur so tief, wie Kapitel 4 ihn braucht (Mut zur Lücke).
%   Keine Formeln außer höchstens der Cosinus-Formel.
% - Reihenfolge so, dass jeder Begriff auf dem vorigen aufbaut:
%   Einbettung -> Cosinus-Ähnlichkeit -> Vektorindex -> lexikalische Suche/BM25
%   -> Rangfusion.
% - Zitate stehen dann HIER, nicht mehr in 4.4 (TODO.md, Block Kapitel 2).
%
% ===== DIE BEGRIFFE (je: was erklären, wo gebraucht, was in Kap. 4 wegfällt) =====
% 1. Einbettung (Embedding)
%    Erklären: Ein Einbettungsmodell wandelt einen Text in einen Vektor fester
%    Länge um. Texte mit ähnlicher Bedeutung erhalten Vektoren mit ähnlicher
%    Richtung. Deutscher Begriff "Einbettung", englisch einmal in Klammern.
%    Gebraucht: 4.3 (Einbettung jedes Knotens aus dem Beschreibungstext),
%    4.4 Stufe 1 ("nur die Aufgabenbeschreibung des Agenten eingebettet"),
%    4.4 semantischer Zweig, 4.4 Stufe 3/4.
%    In Kap. 4 danach: Begriff ohne Erklärung verwenden.
% 2. Cosinus-Ähnlichkeit
%    Erklären: vergleicht die Richtung zweier Vektoren, 1 = gleiche Richtung
%    und damit sehr ähnlicher Inhalt, nahe 0 = kein inhaltlicher Zusammenhang.
%    Quelle: Manning, Raghavan und Schütze 2008, Introduction to Information
%    Retrieval, Kap. 6 (Angaben prüfen, siehe TODO).
%    Gebraucht: 4.4 semantischer Zweig (Schnitt mit Delta und Mindestwert),
%    Stufe 3 (Ähnlichkeit jedes erreichten Knotens), Stufe 4 (Ranking).
%    In Kap. 4 danach streichen: im Absatz "Der erste, semantische Zweig ..."
%    die Sätze "Als Maß dient die Cosinus-Ähnlichkeit [KLÄREN ...]. Sie
%    vergleicht ... Ein Wert von 1 bedeutet ..." -> ersetzen durch kurzen
%    Verweis, z.B. "Maß ist die Cosinus-Ähnlichkeit
%    (Abschnitt~\ref{sec:suche-grundlagen})".
% 3. Vektorindex
%    Erklären: Datenstruktur, die zu jedem Eintrag dessen Vektor speichert und
%    zu einem Anfragevektor schnell die ähnlichsten Einträge findet
%    (nächste Nachbarn). Keine Details zu Verfahren wie HNSW.
%    Gebraucht: 4.3 (je ein Vektorindex für File, Function, Class; Absatz zur
%    Datenbank je Projekt), 4.4 semantischer Zweig.
%    In Kap. 4 danach streichen: Satz "Ein Vektorindex speichert zu jedem
%    Knoten dessen Vektor und findet schnell die Knoten, ..." im Absatz
%    "Der erste, semantische Zweig ...".
% 4. Lexikalische Suche, Volltextindex, BM25
%    Erklären: Suche nach Wörtern statt nach Bedeutung. Ein Volltextindex
%    verzeichnet, welche Wörter in welchen Einträgen vorkommen, und erlaubt
%    auch Präfixanfragen (Einträge, die mit einem Wort beginnen). BM25 ordnet
%    Treffer nach Stichworten und gewichtet seltene Begriffe höher als häufige.
%    Kontrast zur Einbettung: lexikalisch trifft exakte Namen, semantisch trifft
%    Umschreibungen. Genau deshalb kombiniert das Werkzeug beide.
%    Quelle BM25: Robertson und Zaragoza 2009 (siehe TODO).
%    Gebraucht: 4.3 (Volltextindex über Name, Pfad, Beschreibung), 4.4
%    lexikalischer Zweig (Präfixanfrage) und BM25-Zweig.
%    In Kap. 4 danach streichen: im Absatz "Der dritte, lexikalische Zweig ..."
%    den Satz "Sie sucht nach Einträgen, die mit dem Bezeichner beginnen.";
%    im Absatz "Der vierte Zweig, im Folgenden BM25-Zweig, ..." den
%    Nebensatz "ein stichwortbasiertes Rangverfahren, das seltene Begriffe
%    höher gewichtet als häufige [KLÄREN ...]" -> Verweis.
% 5. Rangfusion (Reciprocal Rank Fusion)
%    Erklären: Mehrere Trefferlisten, deren Werte nicht vergleichbar sind,
%    werden über die Plätze zusammengeführt. Jeder Eintrag bekommt aus jeder
%    Liste Punkte nach seinem Platz, die Punkte werden addiert. Wer in mehreren
%    Listen weit oben steht, gewinnt.
%    Quelle: Cormack, Clarke und Buettcher 2009 (siehe TODO).
%    Hier allgemein mit Konstante k: Punkte = 1/(k + r), r = Rangplatz.
%    Gebraucht: 4.4 Absatz "Nach den vier Zweigen liegen vier Listen vor ..."
%    und der folgende Absatz zur Sonderregel des exakten Zweigs.
%    In Kap. 4 danach: den Satz "Vergleichen lässt sich aber der Rangplatz ..."
%    und die allgemeine Erklärung durch Verweis ersetzen. BLEIBEN muss in 4.4:
%    warum die Werte nicht addierbar sind (Ähnlichkeit 0-1, BM25 offen, exakte
%    Liste ohne Werte), die Formel mit dem konkreten Wert k = 60 und das
%    Rechenbeispiel (2/63 gegen 1/61), der ganze Absatz zur Sonderregel des
%    exakten Zweigs (Ersatz-Plätze 1/3/6/20, Aufschlag für häufige Namen,
%    Gewicht 1,2). Das ist Werkzeugentscheidung. Konkrete Werte bleiben im Text
%    (Stilregel seit 30.10.).
%
% ===== NICHT IN DIESEN ABSCHNITT =====
% - Graph, Knoten, Kante, Hop, Seeds, Traversierung: werkzeugspezifisch, werden
%   in 4.1 (sec:werkzeug-ueberblick) eingeführt.
% - tree-sitter, AST, Parser: der Halbsatz in 4.3 reicht.
% - Neo4j, Ollama, Modellname, Parameterwerte: gehören zu Kapitel 4.
%
% ===== BUDGET =====
% Kapitel 2 hat 4 S., 2.1 und 2.2 sind mit je 2 S. verplant. Die 0,75 S.
% müssen dort eingespart werden oder das Kapitelbudget steigt. Mit dem Nutzer
% klären, bevor geschrieben wird. Ersparnis in 4.4 etwa 150-200 Wörter.
%
% ===== NACH DEM SCHREIBEN =====
% 1. In 4.3/4.4 die oben unter "In Kap. 4 danach" genannten Stellen kürzen und
%    Verweise auf sec:suche-grundlagen setzen. Stellen über den zitierten
%    Satzanfang suchen, nicht über Zeilennummern (die verschieben sich).
% 2. KLÄREN-Marker für Cosinus, BM25 und RRF aus 4.4 entfernen, die Zitate
%    stehen dann hier.
% 3. TODO.md: beide Punkte im Block "Kapitel 2" abhaken.
\label{sec:suche-grundlagen}
